\documentclass[12pt,a4paper]{article}
\usepackage{minted}

\makeatletter
\def\input@path{{/home/maenwe/Master_CSM/modèle_cours_latex}}
\makeatother

\usepackage{prise_de_note}


\begin{document}
	
	
	\begin{titlepage}
		\centering
		
		% Logo
		\begin{figure}[h]
			\centering
			% Ajustez la largeur si besoin (ex: 0.4 ou 0.6)
			
\includegraphics[width=0.5\textwidth]{logo.jpg} 
		\end{figure}
		
		\vspace{2cm}
		
		% Filet horizontal du haut
		\rule{\linewidth}{0.5mm} \\[0.4cm]
		
		% Titre
		{\color{bleuTitre} \Huge \bfseries Traveaux Pratiques} \\[0.4cm]
		
		% Filet horizontal du bas
		\rule{\linewidth}{0.5mm}
		
		\vspace{1.5cm}
		
		{\Large \textbf{Éléments finis}}
		
		\vfill
		
		% Bloc Informations
		\begin{minipage}{0.45\textwidth}
			\begin{flushleft} \large
				\textbf{Enseignant :} \\
				M. Eric \textsc{DARRIGAND}
			\end{flushleft}
		\end{minipage}
		\hfill
		\begin{minipage}{0.45\textwidth}
			\begin{flushright} \large
				\textbf{Étudiant :} \\
				Nathan \textsc{HASCOET}
			\end{flushright}
		\end{minipage}
		
		\vspace{2cm}
		
		% Date
		{\large \today}
		
	\end{titlepage}
	
	\section{TP séance 1}
		
	\vspace{-.5cm}
	
	\includegraphics[scale=.6]{photo1.png}
	
	\hspace{-1cm}\textbf{\underline{On construit le maillage (ici des rectangles) :}}
		
	\includegraphics[scale=.7]{photo2.png}
	
	\textit{,pour des triangles, il suffit de rajouter les diagonales, voir sujet TP 1}\\
	
	\textbf{\underline{Maillage :}}
	\begin{enumerate}
		\item $n$ : nombre de nœuds
		\item $x_1$ $y_1$ : coordonnées du nœud 1
		\item \begin{minipage}[t]{10cm}
			$x_2$ $y_2$ : coordonnées du nœud 2
		
			$\displaystyle\vdots$
			
			 $x_n$ $y_n$ : coordonnées du nœud $n$
		\end{minipage}
		\item \begin{minipage}[t]{10cm}
			$m$ : nombre d'éléments (triangles ou quadrilatères ici) \\
			$t$ : type des éléments\\
			$p$ : nombre de nœuds par éléments
		\end{minipage}
		\item \begin{minipage}[t]{15cm}
			$s_1^{(1)}, s_2^{(1)},\ldots,s_p^{(1)}$ : liste des numéraux globaux des nœuds de l'élément 1.
		
		$\displaystyle\vdots$
		
		$s_1^{(m)}, s_2^{(m)},\ldots,s_p^{(m)}$ : liste des numéraux globaux des nœuds de l'élément $m$.
		\end{minipage}
	\end{enumerate}
	
	\vspace{.5cm}
	
	\begin{minipage}{5cm}
		\includegraphics[scale=.7]{photo3.png}
	\end{minipage} \begin{minipage}{10cm}
		\underline{6ème élément :} 5 7 8
	\end{minipage}
	
	\vspace{.5cm}
	
	La \textbf{convention} que l'on adoptera pour la numérotation est la suivante : 
	
	\begin{center}
		\includegraphics[scale=1.2]{photo4.png}
		
		\textit{le numéro local \fbox{3} est celui de l'\underline{angle droit}}\\
		et\\
		\textit{et ensuite nous \underline{tournons en sens direct}}
	\end{center}
	
	\newpage
	
	\subsection{Production perso}	
	
	\vspace{-.5cm}
	
	\begin{minted}{c}
		#include <stdio.h>
		#include <stdlib.h>
		
		void Creer_maillage(double a, double b, double c, double d, int n1, int n2, int t, FILE *fp)
		{
			/*
			===== Cette ALGORITHME Créer un maillage sur un quadrangle =====
			Omega = [a;b] x [c; d] avec a < b  et c < d
			*/
			
			
			int n = n1 * n2; // n = est le nombre de noeuds
			fprintf(fp,"%d\n",n);
			
			double h1 = (b-a)/(n1-1);
			double h2 = (d-c)/(n2-1);
			
			double x,y;
			// Boucle sur les lignes (ordonnée)
			for (int i_y = 0; i_y < n1; i_y++) 
			{
				// Boucle sur les colonnes (abscisse)
				for (int i_x = 0; i_x< n2; i_x++) 
				{
					x = a + h1 * i_x;
					y = b + h2 * i_y;
					fprintf(fp,"(%lf,%lf)\n",x,y);
				}
			}
			
			int m = (n1 - 1)*(n2 - 1); // Nombre d'éléments (triangle ou quadrangle)
			int p;
			switch (t) {
				case 1 : // t = 1 : quadrangle
					p = 4; // nombre de noeuds par éléments
					fprintf(fp,"%d %d %d\n",m, t, p );
					for (int i = 0; i < (n1-2); i++) // boucle sur les lignes
					{
						for (int j = 0; j< (n2-2); j++) // boucle sur les colonnes
						{
							//On récupère le numéro global du 1er sommet de l'élément (i+1)*(j+1)
							int e0 = 2 + n1 * i + j ; 
							fprintf(fp,"%d %d %d %d\n",e0, e0 + n1, e0 + n1 - 1, e0 - 1);
						}
					}
				
				case 2 : // t = 2: triangle
					m = 2*m;
					p = 3; 
					fprintf(fp,"%d %d %d\n",m, t, p );
					for (int i = 0; i < (n1-2); i ++) 
					{
						for (int j = 0; j< (n2-2); j++) 
						{
							 //On récupère le numéro global du 1er sommet du triangle du bas
							int e0 = 2 + n1 * i + j ;
							// Triangle du bas
							fprintf(fp,"%d %d %d\n",e0, e0 + n1, e0 - 1); 
							// Triangle du haut
							fprintf(fp,"%d %d %d\n",e0 + n1 - 1, e0, e0 + n1);
						}
					}
			}
		}
	\end{minted}
	
\section{TP séance 2}
	
	Création du maillage :
	
	\includegraphics[scale=.5]{photo5.png}
	
	$\star$ Ecriture des coordonnées: 
	
	\enc[2cm]{red}{$\begin{matrix}
			x_1&y_1\\
			\vdots & \\
			x_n & v_n
		\end{matrix}$} \rouge{Coordonnée}\\
	
	\rouge{$\rightarrow$ Appel de \textbf{etiqAr} (calcul de \textbf{nRefAr})}
	
	$\star$ Écriture des éléments :
	
	$\underbrace{\enc[3.5cm]{red}{$\begin{matrix}
			s_1^(1)& \cdots&s_p^{(1)} \\
			\vdots & & \\
			s_1^{(m)} & \cdots & s_p^{(m)}
		\end{matrix}$}}_{\text{\rouge{ngnel}}}$ \hspace{.1cm} $\underbrace{\begin{matrix}
		r_1^{(1)} & \cdots & r_q^{(1)} \\
		\vdots & & \\
		r_1^{(m)} &  \cdots & r_q^{(m)}
	\end{matrix}}_{\text{\rouge{nRefAr}}}$\\

	\begin{minipage}{5cm}
		$\star$ Numérotation locale :
	\end{minipage} \begin{minipage}{3cm}
		\includegraphics[scale=.5]{photo6.png}
	\end{minipage}
	
	\underline{Lecture du maillage}\\
	Ecrire une fonction \textit{main} :
	\begin{enumerate}[label=$\to$]
		\item déclare des pointeurs : $\underbrace{\text{coord.}}_{\text{float **}}$, $\underbrace{\text{ngnel, nRefAr}}_{\text{int **}}$
		
		\item appel de lecfima ($\ldots$,\&coord,$\ldots$) $\to$ allocation dynamique de coordonnée.
		
		\item écrire le maillage lu dans un autre fichier
	\end{enumerate}
	
	\subsection{Production Perso}
	
	\begin{minted}{c}
		#include <stdio.h>
		#include <stdlib.h>
		#include "header.h"
		
		void Creer_maillage(double a, double b, double c, double d, int n1, int n2, int t, FILE *fp)
		{
			//====================== Cette ALGORITHME Créer un maillage sur un quadrangle ==============================
			// Omega = [a;b] x [c; d] avec a < b  et c < d
			// n1 = nombre de noeuds en abscisse
			// n2 = nombre de noeuds en ordonnée
			// t = type d'élément : t = 1 pour des quadrangles, t = 2 pour des triangles
			// fp = fichier de sortie
			
			
			int n = n1 * n2; // n = est le nombre de noeuds
			fprintf(fp,"%d\n",n);
			
			double h1 = (b-a)/(n1-1); // 
			double h2 = (d-c)/(n2-1);
			
			double x,y;
			for (int i_y = 0; i_y < n2; i_y++) // Boucle sur les lignes (ordonnée)
			{
				for (int i_x = 0; i_x< n1; i_x++) // Boucle sur les colonnes (abscisse)
				{
					x = a + h1 * i_x;
					y = c + h2 * i_y;
					fprintf(fp,"%lf %lf\n",x,y);
				}
			}
			
			int m = (n1 - 1)*(n2 - 1); // Nombre d'éléments (triangle ou quadrangle)  
			int *p_nbtel = &m; //nbtel = nombre d'éléments 
			// REDONDANT, donné par m
			int p; // p = nombre de noeuds par élément
			int q; // q = nombre d'arêtes par élément
			// REDONDANT, donné par p, car c'est une figure fermée
			//On initialise le tableau des arrêtes
			int *p_nbaret = &p; // nbaret = nombre d'arrête dans un élément 
			// REDONDANT, donné par p
			int nrefdom = 0; // l'intérieur du domaine Omega est numéroté : 0
			int *p_nRefAr, **nRefAr;
			const int nrefcot[] = {1,2,3,4}; // numérotation des bords du domaine Omega
			
			switch (t) {
				case 1 : // t = 1 : quadrangle
				p = 4; // p = 4 (nombre de noeuds par éléments)
				q = *p_nbaret;
				//On initialise le tableau contenant le numéro des arrêtes à 0
				p_nRefAr = malloc(*p_nbaret**p_nbtel*sizeof(*p_nRefAr));
				nRefAr = malloc(*p_nbtel*sizeof(*nRefAr));
				for (int i=0;i<*p_nbtel;i++) nRefAr[i] = p_nRefAr+i**p_nbaret;
				for (int i=0;i<*p_nbtel;i++) {
					for (int j=0;j<*p_nbaret;j++) {
						nRefAr[i][j] = nrefdom;
					}
				}
				
				etiqAr(t, n1, n2, nrefdom, nrefcot, *p_nbtel, *p_nbaret, nRefAr);
				
				fprintf(fp,"%d %d %d %d\n",m, t, p, q);
				for (int i = 0; i < (n1-1); i++) // boucle sur les colonnes
				{
					for (int j = 0; j< (n2-1); j++) // boucle sur les lignes
					{
						int e0 = 2 + n1 * i + j ; 
						//On récupère le numéro global du 1er sommet de l'élément (i+1)*(j+1)
						int numero_element = j + i*(n1-1); // Numéro de l'élément sur lequel on est
						// e0 : sommet en bas à droite ============ e0+ n1 : sommet en haut à droite
						//e0 + n1 -1 : sommet en haut à gauche ==== e0 - 1 : sommet en bas à gauche
						fprintf(fp,"%d %d %d %d %d %d %d %d\n",e0 + n1, e0 + n1 - 1, e0 - 1,
																	nRefAr[numero_element][0],
																	nRefAr[numero_element][1],
																	nRefAr[numero_element][2],
																	nRefAr[numero_element][3]);
					}
				}
				break;
				
				case 2 : // t = 2: triangle
				m = 2*m;
				p = 3; // p = 3
				q = *p_nbaret;
				
				//On initialise le tableau contenant le numéro des arrêtes à 0
				p_nRefAr = malloc(*p_nbaret**p_nbtel*sizeof(*p_nRefAr));
				nRefAr = malloc(*p_nbtel*sizeof(*nRefAr));
				for (int i=0;i<*p_nbtel;i++) nRefAr[i] = p_nRefAr+i**p_nbaret;
				for (int i=0;i<*p_nbtel;i++) {
					for (int j=0;j<*p_nbaret;j++) {
						nRefAr[i][j] = nrefdom;
					}
				}
				
				etiqAr(t, n1, n2, nrefdom, nrefcot, *p_nbtel, *p_nbaret, nRefAr);
				
				fprintf(fp,"%d %d %d %d\n",m, t, p, q);
				for (int i = 0; i < (n1-1); i++) // boucle sur les colonnes
				{
					for (int j = 0; j< (n2-1); j++) // boucle sur les lignes
					{
						int e0 = 2 + n1 * i + j ; //On récupère le numéro global du 1er sommet du triangle du bas
						int numero_element = 2*(j + i*(n1-1)); // Numéro du triangle du bas
						fprintf(fp,"%d %d %d %d %d %d\n", e0, e0 + n1 - 1, e0 - 1,
																				 nRefAr[numero_element][0],
																				 nRefAr[numero_element][1],
																				 nRefAr[numero_element][2]); 
						// Triangle du bas
						fprintf(fp,"%d %d %d %d %d %d\n", e0 + n1 - 1, e0, e0 + n1,
																				 nRefAr[numero_element+1][0],
																				 nRefAr[numero_element+1][1],
																				 nRefAr[numero_element+1][2]);
					// Triangle du haut
					}
				}
				break;
			}
		}
		
	\end{minted}
	
	\newpage
	
	\begin{minted}{c}
		void etiqAr(int type1, int n1, int n2,int nrefdom, const int *nrefcot, int nbtel, int nbaret, int **nRefAr)
		{
			/*
			n1 : nombre de sommets sur une ligne
			n2 : nombre de sommets sur une colonne
			
			type1 = 1 si les éléments du maillage sont des triangles et 2 si Carrés
			nrefdom = intérieur du domaine
			nrefcot = numéro des bords du maillage
			
			nRefAr = tableau de nbtel*nbaret lignes et 3 ou 4 colonnes (suivant si c'est un triangle ou un carré)
			
			VARIABLES REDONDANTES !! 
			nbtel = nombre d'éléments dans le maillage (dépend uniquement du type et de n1 et n2)
			nbaret = nombre d’arête dans un élément (dépend uniquement du type)
			*/
			
			int c = (n1-1)*(n2-1)-(n1-2); // Numéro de l'élément en haut à gauche du maillage
			
			switch (type1) {
				//===================== CARRE : énumération des arrêtes =======================
				
				case 1 :
				//======On numérote les éléments du bord========
				//Ligne du bas du maillage
				for (int i = 0;i<n1-1;i++)
				{
					nRefAr[i][3] = nrefcot[0];
				}
				
				//Colonne de droite
				for (int i = n1-2;i<(n1-1)*(n2-1);i += n1-1)
				{
					nRefAr[i][0] = nrefcot[1];
				}
				
				//Ligne du haut du maillage
				for (int i=c-1;i<(n1-1)*(n2-1);i++)
				{
					nRefAr[i][1] = nrefcot[2];
				}
				
				//Colonne de gauche
				for (int i=0;i<c;i+=n1-1)
				{
					nRefAr[i][2] = nrefcot[3];
				}
				break;
				
				case 2 :
				//===================== TRIANGLE : énumération des arrêtes =======================
				//======On numérote les éléments du bord========
				//Ligne du bas du maillage (on saute un triangle sur 2)
				for (int i = 0;i<2*(n1-1)-1;i+=2)
				{
					nRefAr[i][2] = nrefcot[0];
				}
				
				//Colonne de droite
				for (int i = 2*(n1-1)-1;i<2*(n1-1)*(n2-1);i += 2*(n1-1))
				{
					nRefAr[i][1] = nrefcot[1];
				}
				
				//Ligne du haut du maillage
				for (int i=2*c-1;i<2*(n1-1)*(n2-1);i+=2)
				{
					nRefAr[i][2] = nrefcot[2];
				}
				
				//Colonne de gauche
				for (int i=0;i<2*c-1;i+=2*(n1-1))
				{
					nRefAr[i][1] = nrefcot[3];
				}
				break;
			}
		}
		
	\end{minted}
	
	\newpage
	
	\section{TP séance 3}
	
	$\displaystyle\int_K \nabla w_j^K \cdot \nabla w_i^K$; $\displaystyle\int_K w_j^K \cdot w_i^K$; $\displaystyle\int_K g_N \cdot w_i^K$; $\displaystyle\int_K f_\Omega \cdot w_i^K$\\
	
	\begin{minipage}{12cm}
		Soit $K'$ une arête d'un triangle, $K$ dans $\mathcal{T}_h$ tel que $K' \subset \Gamma_{N,h}$.
	\end{minipage}\begin{minipage}{5cm}
	\includegraphics[scale=.5]{photo7.png}
	\end{minipage}
	
	\[\begin{aligned}
		\displaystyle\int_K a_{0,0} w_j^K w_i^K 
		&= \int_{\widehat K} a_{0,0}(F_K(\hat x))\underbrace{w_j^K(F_K(\hat x))}_{=\hat w_j(\hat x)} \underbrace{w_i^K(F_K(\hat{x}))}_{\hat w_i(\hat x)}\abs{det JF_K(\hat x)}d\hat x\\
		&\simeq \sum_{k=1}^q \hat\omega_k a_{0,0}(F_K(\hat x_k))\hat\omega_j(\hat x_k)\hat\omega_i(\hat w_k)\abs{\det JF_K(\hat x_k)}
	\end{aligned}\]
	
	Besoin de :
	\begin{enumerate}[label=$\star$]
		\item $\displaystyle\sum_k$, $\hat\omega_k$, $\hat x_k$ poids et points d'une formule de quadrature sur $\widehat{K}$.
		
		\item $\hat w_i$, $i=1,\cdots,N^K$ sur $\widehat{K}$, \rouge{et leurs dérivées premières.}
		
		\item $F_K$ définie à partir des $\hat w_i$, $F_K(\hat x) = \displaystyle\sum_{i=1}^{N^K}w_i(\hat x)a_i^K$.\\
		$JF_K$ s'obtient en dérivant l'expression précédente, et on en déduit \rouge{$JF_K^{-1}$}
	\end{enumerate}
\end{document}